\partopening{Theory and the Limits of Identification}{part:theory}
{When does a strategy inherit a bubble, and what can observed prices determine?}

\section{Constant exposure and the two endpoints}\label{sec:classification}
Consider a discounted asset satisfying
\begin{equation}\label{eq:localvol}
 \mathrm{d}S_t=\sigma(S_t)\,\mathrm{d}W_t,\qquad S_0=s_0>0,
\end{equation}
under $Q$. The coefficient is Borel and nowhere zero on $(0,\infty)$, with
$1/\sigma^2$ integrable on every compact subinterval. The nonnegative solution
is absorbed at zero. These assumptions allow weak diffusion laws; differentiability
of the coefficient is unnecessary for the criterion below.

An investor maintains exposure $\pi>0$ to the asset and the rest of wealth in
cash. Before the asset reaches zero, self-financing gives
$\mathrm{d}V_t^\pi=\pi V_t^\pi\mathrm{d}S_t/S_t$. The wealth under study is
defined at all times by
\begin{equation}\label{eq:wealth}
 V_t^\pi=v_0\mathcal E\!\left(\pi\int_0^\cdot
                   \frac{\mathrm{d}S_u}{S_u}\right)_t .
\end{equation}
The exponential is continued at exit according to the convention in the
martingale criterion of Mijatovi\'c and Urusov
\cite{mijatovic2010martingale}: it is set to zero when its quadratic-variation
integral diverges at exit. We call this the \emph{wealth representation}.
It is the specified model of wealth, not a proof that every implementable
strategy through default has this value. In this driftless diffusion the
integral diverges when $S$ reaches zero, so this wealth is then zero.

For any fixed $c>0$, define
\begin{equation}\label{eq:integrals}
 I_0=\int_0^c\frac{x}{\sigma^2(x)}\dd x,
 \qquad I_\infty=\int_c^\infty\frac{x}{\sigma^2(x)}\dd x.
\end{equation}
Local integrability makes their finiteness independent of $c$. The established
driftless criterion says that $S$ is strict on every positive finite horizon
exactly when $I_\infty<\infty$; the corresponding lower-endpoint test says that
zero is accessible exactly when $I_0<\infty$
\cite{mijatovic2010martingale,karatzas2015distribution}.

\begin{theorem}[Constant-proportion classification]\label{thm:exposure}
Under \eqref{eq:localvol}, the stated coefficient assumptions, and the wealth
representation \eqref{eq:wealth}, for every finite $T>0$,
\begin{equation}\label{eq:classification}
 V^\pi\text{ is strict on }[0,T]\quad\Longleftrightarrow\quad
 \left(0<\pi<\tfrac12\ \text{and}\ I_0<\infty\right)
 \ \text{or}\ 
 \left(\pi>\tfrac12\ \text{and}\ I_\infty<\infty\right).
\end{equation}
In particular, $\pi=1/2$ always gives a martingale on each finite horizon.
\end{theorem}

\noindent The result is conditional on the cited martingale criterion and
the wealth representation. Its deterministic endpoint calculation is the
project's proof. The auxiliary-diffusion family already appears in
Karatzas and Ruf, equation (7.4) \cite{karatzas2015distribution}; the
contribution here is the classification across exposure and its implications
for rebalancing and default.

\begin{table}[htbp]\centering
\begin{tabular}{@{}lll@{}}\toprule
Exposure & True martingale on every finite horizon iff & Endpoint that matters\\
\midrule
$0<\pi<1/2$ & $I_0=\infty$ & Zero\\
$\pi=1/2$ & Always & Neither permits auxiliary exit\\
$\pi>1/2$ & $I_\infty=\infty$ & Infinity\\
\bottomrule\end{tabular}
\caption{The exposure classification under the stated wealth convention.}
\end{table}

\subsection{Why one-half appears}
The criterion examines an auxiliary diffusion with drift
$\pi\sigma^2(x)/x$ and the original diffusion coefficient. Its scale density is
$(x/c)^{-2\pi}$. Above one-half, the upper endpoint of its scale is finite;
below one-half, the lower endpoint is finite. The corresponding Feller
integrands reduce to
\begin{equation}\label{eq:fellerkernels}
 \frac{1}{2\pi-1}\frac{x}{\sigma^2(x)}\quad(\pi>1/2),
 \qquad
 \frac{1}{1-2\pi}\frac{x}{\sigma^2(x)}\quad(0<\pi<1/2).
\end{equation}
At one-half the scale is logarithmic and neither endpoint allows this exit.
Neither endpoint is ``good'' in the terminology of the cited criterion, so
auxiliary exit decides strictness. Appendix~\ref{app:proofs} gives the
calculation in detail.

The lower endpoint changes the usual inheritance intuition. Below one-half,
continuous rebalancing into an asset that can reach zero creates strict wealth
even when the asset itself is a martingale. The investor buys more units as
the price declines, continually replenishing the risky fraction from cash.
The model rides this process all the way to zero wealth. That pathwise
description is one part of the exposure classification, which also accounts
for gains on surviving paths through their expectation under $Q$.

\begin{example}[Two reversals at a quarter allocation]\label{ex:reversal}
For absorbed Brownian motion, $\sigma(x)=1$, $I_0<\infty$ and
$I_\infty=\infty$. The asset is a martingale but $V^{1/4}$ is strict.
For $\sigma(x)=x^2$, $I_0=\infty$ and $I_\infty<\infty$.
The asset is strict but $V^{1/4}$ is a martingale. Part~\ref{part:funds}
shows that every finite rebalancing grid reverses the continuous-wealth
classification in both examples.
\end{example}

At $\sigma(x)=1+x^2$, both integrals are finite. Every positive exposure except
one-half then produces strict wealth. The boundary behavior, as well as
the exposure, determines the outcome.
\source{claim 001; experiments 001 and 006}

\subsection{Long horizons, changing exposure, and independent assets}
Three extensions help delimit the classification. Applying the
infinite-horizon criterion of Mijatovi\'c and Urusov, constant nonnegative
exposure gives a \emph{uniformly integrable martingale over all time} only at
$\pi=0$ \cite{mijatovic2010martingale}. Here uniform integrability is required together with the martingale
property. For example, geometric Brownian wealth at positive constant
exposure is a martingale on each finite horizon but fails uniform
integrability over all time.

Changing exposure can restore uniform integrability. For geometric Brownian
motion $\sigma(x)=\kappa x$, with $\kappa>0$, a bounded deterministic schedule
$\pi(t)$ gives a uniformly integrable martingale exactly when
$\int_0^\infty\pi(t)^2\dd t<\infty$. For example,
$\pi(t)=(1+t)^{-\gamma}$ qualifies exactly when $\gamma>1/2$.
For an adaptive schedule, the relevant almost-sure energy condition is under
the auxiliary measure in Ruf's criterion, not simply under $Q$
\cite{ruf2015martingale}. Bounded exposure alone is insufficient for a general
local-volatility asset.

If two full asset--wealth systems are independent under $Q$ and their
product is a continuous local martingale in the common filtration, the
product wealth has multiplying mean factors. Writing their normalized defects as
$\delta_1$ and $\delta_2$, the product defect is
\begin{equation}\label{eq:productdefect}
 \delta_{12}=\delta_1+\delta_2-\delta_1\delta_2.
\end{equation}
It vanishes exactly when both defects vanish. An independent
true-martingale factor does not remove the strictness of another factor.
Independence of the full systems is the key assumption behind this
factorization.
\source{claims 074--076; numerical checks summarized in the
original paper's Section 3.1}

\section{What prices identify}\label{sec:identification}
The classification depends on what happens near zero and infinity. A finite
record occupies a finite price range. Better measurement inside that range
does not by itself supply the missing behavior outside it. This is the
extrapolation problem already recognized by Jarrow, Kchia, and Protter
\cite{jarrow2011detect}.

\begin{proposition}[Compatible tails]\label{prop:tails}
Suppose $\sigma$ is known exactly on $(0,R]$, with $c<R$. Set
$J=\int_c^R x/\sigma^2(x)\dd x$. Over all admissible coefficients agreeing
with $\sigma$ on that interval, the upper-endpoint integral ranges over
$(J,\infty]$. Its finiteness can take either value. If $\sigma$ is positive
and continuous, the compatible coefficients can be chosen positive and
continuous. If only $[r,R]$, $r>0$, is known, the lower and upper unknown tails
can be varied independently.
\end{proposition}

\begin{proof}
Keep $\sigma$ unchanged below $R$ and set
$g_p(x)=\sigma(R)(x/R)^p$ above it. For $p=1$ the tail integral diverges;
for $p>1$ it equals
\[
 \int_R^\infty \frac{x}{g_p(x)^2}\dd x
 =\frac{R^2}{2(p-1)\sigma(R)^2}.
\]
Varying $p$ yields every strictly positive finite tail contribution, together
with infinity. A similar completion below $r$ varies the lower tail. The
known middle contribution is unaffected.
\end{proof}

This proposition grants much more than actual observations: exact knowledge
of a coefficient on an interval. It is an identified-set statement, not a
confidence interval or a finite-sample testing theorem. The latter requires
comparing observation laws.

\begin{theorem}[Finite-horizon testing obstruction]\label{thm:nogo}
Fix a positive finite horizon and finitely many observed price paths. In the
class of driftless local-volatility diffusions described above, including
$P=Q$, every bounded randomized test has power at each strict alternative no
larger than its rejection probability maximized over all true-martingale
nulls. The conclusion persists when coefficients are required to be positive
and continuous. It also has a version allowing equivalent physical measures.
\end{theorem}

For $P=Q$, the argument is particularly transparent. Replace the coefficient
above a large level $n$ by a linear tail. The completed price is a martingale
and agrees in stopped law with the alternative until the price reaches $n$.
The probability of reaching that level before the horizon tends to zero, so
the observed path laws become arbitrarily close. A finite number of paths
does not change the argument. Weak existence and stopped uniqueness in law
are cited diffusion inputs; their full reconstruction is not part of the
project. Appendix~\ref{app:proofs} records the probability comparison.

The theorem concerns \emph{uniform} size over a broad null class.
Restricted models can support informative tests by excluding the compatible
completions, for example through a specified tail law. Under $P=Q$, a known
full law identifies $\E_Q[S_T]$; the obstruction concerns inference from
finite records uniformly across the admitted laws.
\source{claims 002, 003, 006, and 013}

\subsection{Sin's model: feedback and the pricing measure}\label{sec:sin}
Sin's stochastic-volatility model provides a second theoretical reason to
study positive price--volatility association \cite{sin1995complications}.
Under a specified pricing probability $Q^s$, write its dynamics as
\begin{equation}\label{eq:sin}
 \mathrm{d}S_t=S_t v_t^\alpha\gamma\cdot\mathrm{d}W_t,\qquad
 \mathrm{d}v_t=v_t a\cdot\mathrm{d}W_t+\kappa(L-v_t)\mathrm{d}t,
\end{equation}
where $W$ is two-dimensional Brownian motion, $\alpha,\kappa,L>0$,
$a,\gamma$ are nonzero constant vectors, $S_0>0$, and $v_0=1$.
The price and volatility factor remain positive. Sin's Theorem 3.2 says
that $S$ is a martingale on all finite horizons exactly when
$a\cdot\gamma\leq0$. Positive alignment therefore produces strictness
on some finite horizon.

The dot product is also the sign of the instantaneous association between
price and volatility shocks. When it is positive, upward price shocks
arrive with upward volatility shocks. Under the auxiliary change of measure
used to examine the price's expectation, that alignment adds an upward
term of order $v^{1+\alpha}$ to the volatility drift. The resulting
explosion mechanism creates the expectation shortfall. This is a
stochastic-volatility route to a bubble, alongside the local-volatility
route in which the diffusion coefficient grows rapidly with the price.

That connection changes how a signal should be interpreted. A positive
fitted level exponent may summarize explicit local-level dependence or
correlated volatility shocks. Those are different explanations of the
coefficient, but both warrant attention when studying bubble-related
price--variance behavior. The simulation response to correlation in
Part~\ref{part:signals} is informative about that broader sensitivity.
The model must still be specified to turn the association into a strictness
criterion: Sin's multiplicative volatility dynamics differ from the
square-root variance dynamics used in those simulations.

Sin's model also separates this mechanism from identification of a pricing
measure. With nonparallel vectors,
$|a\cdot\gamma|<|a||\gamma|$, his Theorem 3.6 supplies an equivalent
measure under which the same price is a martingale on any specified finite
horizon. Taking $P=Q^s$ and choosing a horizon with strictness gives one
physical law compatible with both positive and zero pricing defects.
The change of measure alters the volatility drift while preserving the
instantaneous covariance structure. Thus covariance information helps
characterize the observed feedback, while a valuation conclusion additionally
requires a pricing law. The two roles of Sin's model---a positive feedback
mechanism and pricing-measure ambiguity---belong together.
\source{claim 003, part 3; Sin's Theorems 3.2 and 3.6, as recorded in AX-05}

\subsection{What is positively identified inside the range}
There is a useful positive result. For a continuously observed path of the
diffusion, let $L_T^x$ be semimartingale local time, normalized by the
occupation identity
$\int_0^T f(S_t)\sigma^2(S_t)\dd t=\int f(x)L_T^x\dd x$.
With zero contribution when the denominator is zero, the visited-range
identity is
\begin{equation}\label{eq:visited}
 \int_0^T\frac{S_t}{L_T^{S_t}}\,
     \ind_{\{S_t>0,L_T^{S_t}>0\}}\dd t
 =\int_{(\min S,\max S)\cap(0,\infty)}
        \frac{x}{\sigma^2(x)}\dd x.
\end{equation}
It is an application of the occupation formula and positivity of local time
on the visited interior. The identity is pathwise and transfers under
equivalent physical measures. It identifies the visited portion, not its
unobserved completion. Estimating it from sampled prices introduces another
problem, especially near levels with little occupation.

Under a declared common power tail, pooling names expands the union of
visited price levels. Claim 015 gives a limiting detection probability at
least as large as that for any one name, without assuming independence.
This is a coverage gain; no estimation-rate improvement is asserted.
Separate reviewer simulations of the later pooled likelihood constructions
in claims 062--063 found no improvement over per-name testing.
\source{claims 008, 012, 013, and 015; experiment 013 for sampled
visited-integral estimation; reviewer simulations in claims 062--063 for pooling}

\subsection{Why sampling frequency is not a universal cure}
In a separable model for log price $X$, suppose its instantaneous variance is
$v_t=\exp(a(t)+\beta X_t)$ and $a$ is deterministic, continuous, and of finite
variation. Continuous observation yields
\begin{equation}\label{eq:continuousbeta}
 \beta=\frac{[\log v,X]_T}{[X]_T},
 \qquad v_t=\frac{\mathrm{d}[X]_t}{\mathrm{d}t}.
\end{equation}
Finite variation contributes no quadratic covariation, so this follows by
polarization. Stochastic volatility with nonzero quadratic variation does
not satisfy that premise. Nor does \eqref{eq:continuousbeta} provide a
noise-robust estimator of the derivative of quadratic variation from daily
prices.

Under $P=Q$ and the cited exact price laws, on every fixed finite grid the
normalized laws of
$\mathrm{d}S=\kappa S\,\mathrm{d}W$ and
$\mathrm{d}S=\kappa S^2\,\mathrm{d}W$, started at one, approach one another in
total variation as $\kappa\downarrow0$. The first is a martingale and the
second strict, with a fixed separation in elasticity. Thus no uniform
finite-grid power or required-frequency guarantee follows from elasticity
separation and a time-variation budget alone. A lower bound on the relevant
variance scale is missing. At a given scale, richer intraday observations can still provide the
substantial gains documented in the next part.
\source{claims 071--073}
